\subsection{Sluoksnių aprašymas}
\label{subsec:sluoksniai}

Konvoliucinis neuroninis tinklas sudarytas iš kelių skirtingų tipų sluoksnių, kurie atlieka skirtingas funkcijas. Šiame poskyryje aprašoma kiekvieno projekte naudojamo sluoksnio paskirtis ir veikimo principas.

\subsubsection{Konvoliucinis sluoksnis}
\label{subsubsec:konvoliucinis_sluoksnis}

Konvoliucinis sluoksnis (angl.~\textit{convolutional layer}) yra pagrindinis CNN požymių išskyrimo elementas. Jo veikimo principas – nedidelis filtras (branduolys, angl.~\textit{kernel}) stumiamas per įėjimo duomenis ir kiekvienoje pozicijoje apskaičiuojama skaliarinė sandauga tarp filtro svorių ir atitinkamo įėjimo fragmento. Rezultatas – požymių žemėlapis (angl.~\textit{feature map}), rodantis, kur aptiktas filtro ieškomas šablonas.

Projekte naudojami du konvoliucijų tipai:
\begin{itemize}
    \item \texttt{Conv2d} – 2D konvoliucija vaizdams. Filtras stumiamas dviem kryptimis (aukštis ir plotis). Pvz., $3 \times 3$ dydžio branduolys nuskaito vaizdo sritį po $9$ pikselius ir aptinka erdvinius šablonus – kraštus, tekstūras, formas.
    \item \texttt{Conv1d} – 1D konvoliucija laiko eilutėms. Filtras stumiamas viena kryptimi (sekos ašimi). Projekte jis slenka per $31$ klaviatūros laikinius požymius ir aptinka lokalinius laikinus šablonus.
\end{itemize}

Svarbi konvoliucijų savybė – \textbf{svorių dalinimasis} (angl.~\textit{weight sharing}): tas pats filtras taikomas visose įėjimo pozicijose. Dėl to konvoliucinis sluoksnis turi žymiai mažiau parametrų nei pilnai sujungtas sluoksnis ir geba aptikti šablonus nepriklausomai nuo jų vietos įėjimo duomenyse. Kiekvienas filtras sukuria vieną požymių žemėlapį, todėl filtrų skaičius lemia išėjimo kanalų skaičių.

Projekte visuose konvoliuciniuose sluoksniuose naudojamas branduolio dydis $3$ ir \texttt{padding}~$= 1$, todėl erdviniai matmenys po konvoliucijos išlieka nepakitę – juos mažina tik telkimo sluoksniai.

\subsubsection{Aktyvacijos funkcija}
\label{subsubsec:aktyvacijos_funkcija}

Aktyvacijos funkcija (angl.~\textit{activation function}) įveda netiesiškumą po kiekvieno konvoliucinio ar pilnai sujungto sluoksnio. Be aktyvacijos funkcijos kelių tiesinių sluoksnių kompozicija liktų tiesinė transformacija ir tinklas negalėtų išmokti sudėtingų šablonų. Projekte naudojamos keturios aktyvacijos funkcijos:

\begin{itemize}
    \item \textbf{ReLU} (angl.~\textit{Rectified Linear Unit}) – $f(x) = \max(0,\, x)$. Paprasčiausia ir dažniausiai naudojama aktyvacijos funkcija. Neigiamos reikšmės nulinamos, teigiamos paliekamos nepakeistos. Trūkumas – „mirštančių neuronų" problema: jei neuronas visada gauna neigiamą įėjimą, jo gradientas tampa $0$ ir jis nustoja mokytis.
    \item \textbf{LeakyReLU} – $f(x) = \max(\alpha x,\, x)$, kur $\alpha = 0{,}01$. Sprendžia „mirštančių neuronų" problemą, nes neigiamoms reikšmėms suteikiamas mažas nuolydis ($\alpha$), užtikrinantis, kad gradientas niekada netaptų visiškai nuliniu.
    \item \textbf{Tanh} – $f(x) = \tanh(x)$, reikšmės perkeliamos į intervalą $[-1,\, 1]$. Funkcija yra glodi ir simetriška, tačiau kraštuose prisotina (angl.~\textit{saturates}) – gradientai tampa labai maži, kas lėtina mokymąsi.
    \item \textbf{ELU} (angl.~\textit{Exponential Linear Unit}) – $f(x) = x$, kai $x > 0$, ir $f(x) = \alpha (\ee^{x} - 1)$, kai $x \leq 0$. Neigiamoje dalyje funkcija yra glodi, todėl gradientai geriau sklidžia nei su ReLU.
\end{itemize}

Bazinėje konfigūracijoje naudojama \textbf{ReLU}, o likusiomis funkcijomis tiriama aktyvacijos įtaka modelio tikslumui.

\subsubsection{Telkimo sluoksnis}
\label{subsubsec:telkimo_sluoksnis}

Telkimo sluoksnis (angl.~\textit{pooling layer}) mažina požymių žemėlapio erdvinius matmenis, pasilikdamas svarbiausią informaciją. Projekte naudojamas \textbf{maksimalusis telkimas} (angl.~\textit{max pooling}) – iš kiekvieno telkimo lango pasirenkama didžiausia reikšmė:

\begin{itemize}
    \item \texttt{MaxPool2d} – vaizdams. Stumiamas $2 \times 2$ langas, kuris kiekvienoje pozicijoje iš keturių reikšmių pasirenka didžiausią. Po kiekvieno telkimo sluoksnio vaizdo kraštinė sumažėja dvigubai.
    \item \texttt{MaxPool1d} – laiko eilutėms. Stumiamas $2$ dydžio langas, kuris iš dviejų gretimų reikšmių pasirenka didžiausią. Po kiekvieno telkimo sluoksnio sekos ilgis sumažėja dvigubai.
\end{itemize}

Telkimo sluoksnis atlieka tris funkcijas:
\begin{enumerate}
    \item \textbf{Sumažina skaičiavimų apimtį} – mažesni požymių žemėlapiai reiškia mažiau parametrų tolesniuose sluoksniuose.
    \item \textbf{Suteikia transliacijos invariantiškumą} (angl.~\textit{translation invariance}) – nedidelis šablono pasislinkimas neturi didelės įtakos telkimo rezultatui.
    \item \textbf{Didina receptyvų lauką} (angl.~\textit{receptive field}) – kiekvienas neuronas gilesnėse pakopose „mato" didesnę pradinių duomenų sritį.
\end{enumerate}

Projekte visuose telkimo sluoksniuose naudojamas telkimo dydis $2$.

\subsubsection{Paketinis normalizavimas}
\label{subsubsec:paketinis_normalizavimas}

Paketinio normalizavimo sluoksnis (angl.~\textit{batch normalization}) normalizuoja kiekvieno mini paketo (angl.~\textit{mini-batch}) aktyvacijas taip, kad jų vidurkis būtų artimas $0$, o standartinis nuokrypis – artimas $1$. Po normalizavimo taikoma afininė transformacija su mokomais parametrais $\gamma$ (mastelis) ir $\beta$ (poslinkis), leidžiančiais tinklui išlaikyti reikiamą aktyvacijų pasiskirstymą.

Projekte naudojami du paketinio normalizavimo variantai:
\begin{itemize}
    \item \texttt{BatchNorm2d} – vaizdams. Normalizuoja kiekvieną požymių žemėlapį (kanalą) atskirai per visą mini paketą.
    \item \texttt{BatchNorm1d} – laiko eilutėms. Normalizuoja kiekvieną kanalą per sekos matmenį ir mini paketą.
\end{itemize}

Paketinis normalizavimas teikia tris pagrindinius privalumus:
\begin{enumerate}
    \item \textbf{Stabilizuoja mokymą} – mažina vidinį kovariatinį poslinkį (angl.~\textit{internal covariate shift}), todėl tinklas greičiau konverguoja.
    \item \textbf{Leidžia naudoti didesnį mokymosi greitį} – normalizuotos aktyvacijos sumažina gradientų sproginėjimo (angl.~\textit{exploding gradients}) riziką.
    \item \textbf{Švelnus reguliarizacijos efektas} – mini paketo statistikų naudojimas įneša šiek tiek triukšmo, kas veikia panašiai kaip papildoma reguliarizacija.
\end{enumerate}

Bazinėje konfigūracijoje paketinis normalizavimas yra \textbf{išjungtas}, o jo įtaka tiriama atskirame tyrime.

\subsubsection{Dropout sluoksnis}
\label{subsubsec:dropout_sluoksnis}

\textit{Dropout} sluoksnis yra reguliarizacijos technika, kuri mokymo metu atsitiktinai „išjungia" dalį neuronų su tikimybe~$p$, nustatydama jų aktyvacijas lygias nuliui. Tai priverčia tinklą mokytis perteklinių reprezentacijų ir nesusieti klasifikavimo tik su konkrečiais neuronais, todėl sumažėja persimokimo (angl.~\textit{overfitting}) rizika.

Projekte naudojami du \textit{Dropout} variantai:
\begin{itemize}
    \item \texttt{Dropout2d} – vaizdų tinkle (\texttt{ImageCNN}). Užuot išjungęs atskirus neuronus, išjungia visą požymių žemėlapį (kanalą). Tai efektyviau konvoliuciniams sluoksniams, nes gretimi pikseliai yra stipriai koreliuoti.
    \item \texttt{Dropout} – laiko eilučių tinkle (\texttt{KeystrokeCNN}) ir abiejų tinklų klasifikatoriuje. Atsitiktinai išjungia atskirus neuronus.
\end{itemize}

Testavimo ir validavimo metu \textit{Dropout} yra automatiškai išjungiamas (aktyvuojamas tik \texttt{model.train()} režimu), o aktyvacijos yra atitinkamai mastelizuojamos, kad išėjimo reikšmių vidurkis sutaptų su mokymo metu gautu vidurkiu.

Bazinėje konfigūracijoje \textit{Dropout} tikimybė yra $p = 0{,}0$ (t.\,y. \textit{Dropout} neveikia), o skirtingos tikimybės ($0{,}25$; $0{,}5$; $0{,}75$) tiriamos atskirame tyrime.

\subsubsection{Suplojimo ir pilnai sujungti sluoksniai}
\label{subsubsec:suplojimas_ir_pilnai_sujungtas}

\textbf{Suplojimo sluoksnis}. \texttt{Flatten} sluoksnis pertvarko daugiamačius požymių žemėlapius į vienmatį vektorių. Jis yra tiltas tarp konvoliucinės dalies (požymių išskyrimo) ir klasifikatoriaus (pilnai sujungtų sluoksnių). Sluoksnis neturi mokomų parametrų – jis tik pertvarko tensoriaus formą:
\begin{itemize}
    \item Vaizdų tinkle: $(128,\, 8,\, 8) \to (8192)$ – $128$ kanalų po $8 \times 8$ pikselių suplojama į vieną $8192$ ilgio vektorių.
    \item Laiko eilučių tinkle: $(128,\, 7) \to (896)$ – $128$ kanalų po $7$ laiko žingsnius suplojama į vieną $896$ ilgio vektorių.
\end{itemize}

\textbf{Pilnai sujungtas sluoksnis.} Pilnai sujungtas sluoksnis (angl.~\textit{fully connected layer}, \texttt{Linear}) sujungia kiekvieną įėjimo neuroną su kiekvienu išėjimo neuronu per mokomus svorius ir poslinkius. Projekte klasifikatorių sudaro du tokie sluoksniai:
\begin{enumerate}
    \item \textbf{Paslėptasis sluoksnis} – mažina suploto vektoriaus dimensiją iki \texttt{classifier\_hidden\_size} ($256$ vaizdų tinkle, $128$ laiko eilučių tinkle). Po jo taikoma aktyvacijos funkcija ir \textit{Dropout}.
    \item \textbf{Išėjimo sluoksnis} – turi tiek neuronų, kiek yra klasių ($2$ vaizdų tinkle, $30$ laiko eilučių tinkle). Kiekvienas neuronas išveda neapdorotą balą (angl.~\textit{logit}) atitinkamai klasei. Didžiausią balą turinčios klasės indeksas laikomas modelio spėjimu.
\end{enumerate}

Išėjimo sluoksnis nenaudoja aktyvacijos funkcijos, nes \texttt{CrossEntropyLoss} paklaidos funkcija viduje taiko \textit{Softmax} transformaciją, kuri paverčia neapdorotus balus į tikimybių pasiskirstymą.

